\section{Example Models}
\label{sec:examplemodels}

Four example models instantiate the source metamodel, and each is chosen to
exercise something the metamodel claims. They are committed beside the
metamodels, together with the expected result of parsing each one, so that the
examples are checked rather than illustrative.

\begin{table}[!htbp]
\centering
\small
\begin{tabular}{>{\raggedright\arraybackslash}p{0.16\linewidth}>{\raggedright\arraybackslash}p{0.76\linewidth}}
\hline
\textbf{Example} & \textbf{What it exercises} \\
\hline
minimal & The smallest complete model: one project, one application, one
  process, and nothing that is not required. Everything absent from it is
  either derived or refused, which makes it the model that shows what the
  language forces an author to say. \\
auth & Two processes in one application, so the release gate has more than one
  member and the gate's deadline is a maximum over members rather than a single
  value. \\
data & Several applications in one project sharing a namespace, with storage
  and the backup plans their durability classes earn. \\
knowledge & Dependencies across applications, secrets, a volume already bound
  to a machine, and an exposure whose routes carry different audiences. \\
\hline
\end{tabular}
\caption{The Example Models and What Each One Exercises}
\label{tab:examples}
\end{table}

\subsection{The Minimal Model}

The minimal model is the one to read first. It declares a project with one
application and one process, and it is shorter than the file the same deployment
needs today by roughly an order of magnitude. What it does \emph{not} declare is
the point: no rollout strategy, no resource limits, no security context, no
labels, no selector, no namespace, no probe timings, no instance count. Each of
those either derives from something the model does declare, or is platform
policy, or is refused outright.

\subsection{Refused Models}

Alongside the valid examples, a set of deliberately invalid models is
committed, one per constraint, each with the diagnostic it must produce. A
constraint that has only ever been satisfied is not evidence that it fires, so
each constraint is paired with a model that violates it.

The refused set covers, among others: a route naming a surface its process does
not offer; two routes on one hostname claiming the same path; an alert class
with nothing to collect metrics from; a process declaring a cutover its storage
cannot honour; a secret delivery that the cluster's encryption setting forbids;
a durability class with no engine to back it up; and an exposure whose audience
no edge tier serves.

\subsection{Example Target Models}

A hand-written target model accompanies the minimal example: the resolved model
that the Task 2 transformation must produce from it, written and reviewed by
hand rather than generated. It is the expected output of the transformation and
the input the Task 3 templates are first run against, which is what lets both
tasks be checked before either is complete.
